\documentclass[11pt]{article}
\usepackage[a4paper,margin=24mm]{geometry}
\usepackage{amsmath,amssymb,amsthm,hyperref}
\newtheorem{theorem}{Theorem}
\newtheorem{lemma}{Lemma}
\newtheorem{corollary}{Corollary}
\newcommand{\Sig}{\operatorname{Sig}}
\title{Fast positive prefix completion by doubling the fairness order}
\author{Proof with three completed independent mathematical reviews}
\date{30 September 2026}
\begin{document}\maketitle
The only external existence input in this note is the small $t$-wise
permutation theorem of Kuperberg--Lovett--Peled. A known fair filler with
$\exp(O(n\log n))$ faces per letter is also used, for example the
reviewed factorial-squared construction in this research directory.

\begin{theorem}\label{main}
There is an absolute constant $C$ such that every positive word $u$ on
$n\ge2$ letters is a literal prefix of an equal-composition
permutation-fair word whose count on each letter is at most
\[
 (|u|+1)\exp(Cn\log(n+1)).
\]
In particular positive prefix completion has only a factorial-order
multiplicative cost, rather than the earlier
$\exp(O(n^2\log n))$ repeated full-symmetrization cost.
\end{theorem}

\begin{lemma}[Disjoint equal-power-sum sets]\label{moments}
For integers $q\ge2,d\ge0$, there are $q$ disjoint nonempty subsets
$S_1,\ldots,S_q$ of $\{0,\ldots,N-1\}$ with identical sums of powers
through degree $d$, including degree zero, where
\[
 N\le Cq(d+1)^4\log^2(2q(d+1)).
\]
They may be translated so that the minimum of their union is zero.
\end{lemma}
\begin{proof}
For $d=0$, use $q$ singleton sets. Suppose $d\ge1$ and put
\[
 k=\lceil100(d+1)^2\log(2q(d+1))\rceil,\qquad N=4qk^2.
\]
Partition the $k$-subsets of $\{0,\ldots,N-1\}$ by their sums of powers
of degrees $1,\ldots,d$. The number of possible vectors is at most
\[
 B=(2k)^dN^{d(d+1)/2}.
\]
There is a class of more than $k!(q-1)^k$ sets: indeed
$\binom Nk\ge(N/k)^k$, $k!\le k^k$, and $4^k>B$ imply this assertion.
For the last inequality, let $L=\log(2q(d+1))\ge\log8$.
Then $\log k\le6L$, $\log N\le14L$, and
$\log B\le14(d+1)^2L<k\log4$.

The elementary sunflower lemma says that a family of more than
$k!(q-1)^k$ distinct $k$-sets contains $q$ sets with all pairwise
intersections equal to one common core. Here is the standard induction:
if there are $q$ disjoint members, they suffice; otherwise the union
of a maximal disjoint subfamily has at most $k(q-1)$ elements and
meets every set. Some element lies in more than
$(k-1)!(q-1)^{k-1}$ members. Remove it and apply the induction for
$(k-1)$-sets, then restore it to the core. The base case $k=1$ is
immediate.
Remove the core from the resulting $q$ equal-moment sets.
Their petals are disjoint, nonempty, and retain identical sums of
powers, since the same core was removed from each.
Translation preserves every equality of power sums of degree at most
$d$, by the binomial theorem.
\end{proof}

We work in the rational associative algebra on $X_1,\ldots,X_n$ in
which every word with a repeated letter is zero. Let $\mathcal A_{\ge p}$
be its degree filtration. The signature of a positive word is
$\Sig(i_1\cdots i_L)=\prod_a(1+X_{i_a})$; its coefficients count
scattered subwords. Set $H=\sum_iX_i$.
A balanced word with $M$ copies of every letter is fair through degree
$t$ precisely when
$\Sig(w)=\exp(MH)\pmod{\mathcal A_{\ge t+1}}$.

\begin{lemma}[Averaging a balanced signature]\label{average}
Let $w$ have $M$ copies of every letter, and let $T\subset S_n$ be a
$t$-wise uniform family. Then
\[
 \sum_{\sigma\in T}\bigl(\Sig(\sigma w)-\exp(MH)\bigr)
 \in\mathcal A_{\ge t+1}.
\]
\end{lemma}
\begin{proof}
For any ordered $k$-tuple of distinct letters, $k\le t$, the average
coefficient under $T$ equals the average over all ordered distinct
$k$-tuples in $w$. For each fixed $k$-subset, summing over its orders
counts all $M^k$ selections of one occurrence per letter.
Hence that average is
$\binom nkM^k/(n)_k=M^k/k!$, exactly the coefficient of $\exp(MH)$.
Lower-wise uniformity follows from $t$-wise uniformity by summing over
extensions of a tuple.
\end{proof}

\begin{lemma}[One doubling step]\label{double}
Suppose $w$ and a completely fair filler $g_M$ both have $M$ copies of
every letter, and $\Sig(w)-\exp(MH)\in\mathcal A_{\ge p}$, where
$2\le p\le n$. There is a balanced word $w'$ literally starting with
$w$, fair through degree $t=\min(2p-1,n)$, whose common count is
$M N'$ with
\[
 N'\le\exp(Cp\log(n+1)).
\]
\end{lemma}
\begin{proof}
Theorem~1.6 of Kuperberg--Lovett--Peled gives a $t$-wise uniform set
$T\subset S_n$ of size $q\le(cn)^{ct}$, for a universal constant $c$.
Left-translate it so it contains the identity; this preserves uniformity.
Apply Lemma~\ref{moments} with $d=t-p\le p-1$.
Translate the resulting position sets so that their union starts at zero,
and assign the identity permutation to the color containing zero.
At each selected position of color $\sigma$, place the block $\sigma w$;
fill every other position with $g_M$. The resulting word begins with $w$.
The number $N'$ of positions is at most the $N$ in that lemma, and
$\log N=O(p\log(n+1))$.

Write $G=\exp(MH)$ and $D_\sigma=\Sig(\sigma w)-G$.
Each $D_\sigma$ begins in degree $p$. Expanding the block product
through degree $t<2p$, all terms with two defects vanish, leaving
\begin{equation}\label{linear}
 G^{N'}+
 \sum_{\sigma\in T}\sum_{l\in S_\sigma}
               G^lD_\sigma G^{N'-l-1}.
\end{equation}
For a homogeneous piece of $D_\sigma$ of degree $r\ge p$, its
contribution to degree $k\le t$ is a sum over $a+b=k-r$ of
\[
 \frac{M^{a+b}}{a!b!}\,l^a(N'-l-1)^b H^a(D_\sigma)_rH^b.
\]
Every coefficient in $l$ is a polynomial of degree at most
$k-r\le t-p=d$. Thus the identical power sums of the color classes
turn the second term of \eqref{linear}, degree by degree, into a
common linear operator applied to $\sum_\sigma D_\sigma$.
Lemma~\ref{average} makes this zero through degree $t$.
The remaining signature is $G^{N'}=\exp(MN'H)$, proving the claim.
\end{proof}

\begin{proof}[Proof of Theorem~\ref{main}]
Fix one completely fair word $g$ with common count
$F\le\exp(C_0n\log(n+1))$, supplied for instance by
\path{factorial_squared_construction.tex}.
Append letters to $u$ until every count is the same positive multiple
$M$ of $F$, choosing the smallest such multiple at least the largest
original count. Then $M\le |u|+F$. Call the balanced word $w$.
It is automatically fair through degree one, so the first defect
degree is at least $p=2$. Since $F\mid M$, the filler
$g_M=g^{M/F}$ has exactly the same common count.

Apply Lemma~\ref{double}. Every later common count remains a multiple
of $F$, so the same fixed filler, repeated to the current count, is
available at every stage with no new count multiplier.
The least possible defect degree doubles: $p=2,4,8,\ldots$ until it
exceeds $n$. The logarithms of all count multipliers sum to at most
\[
 C\log(n+1)\sum_{2^j\le n}2^j=O(n\log(n+1)).
\]
Every step retains its input word as a literal prefix. The final word
is fully fair, still begins with $u$, and has common count at most
$(|u|+F)\exp(C_1n\log(n+1))$, which has the stated bound after changing
$C$. The empty input word is covered by the same choice $M=F$.
\end{proof}

\paragraph{Primary external input.}
G.~Kuperberg, S.~Lovett, R.~Peled,
\emph{Probabilistic existence of regular combinatorial structures},
Geom. Funct. Anal. \textbf{27} (2017), 919--972, Theorem~1.6.
The primary version \url{https://arxiv.org/abs/1302.4295}, v3, states
for every $n\ge1$ and $1\le t\le n$ the existence of a $t$-wise uniform
set of permutations of size at most $(cn)^{ct}$, with universal $c$.
The result here is a derived completion statement; no claim is made
that the cited design theorem or the elementary sunflower lemma is new.
\end{document}
